%------------------------------------------------------------------------------%
\documentclass[a4paper,12pt,fleqn]{article}

\usepackage{calculo_varias_variaveis-1}

\ShowAnswers

%------------------------------------------------------------------------------%
\begin{document}

\makeHeader{CFVVI}{Prova Substitutiva -- Turma 7}{21/07/2025}

\textbf{O exercício correspondente a prova que será substituída vale 50 pontos}

% 01
%------------------------------------------------------------------------------%
\exercise{25}
Considerando a função
\(
  f(x, y) = \sqrt{2x-y}-2
\)
\begin{enumerate}[label=\alph*)]
  \item Determine e esboce o domínio de $f$
  \item Determine as curvas de nível de $f$ e esboce três delas
  \item Calcule o limite
  \(
    \lim_{(x, y)\to(2, 0)} \frac{f(x, y)}{2x-y-4}
  \)
  ou mostre que ele não existe
\end{enumerate}
\begin{questiononly}
\begin{multicols}{2}
  Domínio
  \vspace{-0.8\baselineskip}
  \begin{center}
    \includegraphics{axis_xy}
  \end{center}
  Curvas de nível
  \vspace{-0.8\baselineskip}
  \begin{center}
    \includegraphics{axis_xy}
  \end{center}
\end{multicols}
\end{questiononly}
\clearpagequestiononly

\begin{answer}
  \begin{multicols}{2}
    Domínio
    \vspace{-0.8\baselineskip}
    \begin{center}
      \includegraphics{T7_dominio}
    \end{center}
    Curvas de nível
    \vspace{-0.8\baselineskip}
    \begin{center}
      \includegraphics{T7_curvas_nivel}
    \end{center}
  \end{multicols}

\clearpage
%\setlength\columnsep{3em}
\begin{multicols*}{2}
\raggedcolumns
\setlength{\jot}{5pt}

\textbf{a)\;}
O domínio de $f$ consiste dos pontos onde é possível avaliar
a raiz quadrada, isto é, os pontos que satisfazem $2x-y\geq 0$,
ou seja,
\begin{align*}
  2x-y & \geq 0  \\
  -y   & \geq 2x \\
  y    & \leq 2x
\end{align*}
Portanto,
\[
  D_f = \left\{(x, y)\in\R^2 \st y \leq 2x \right\},
\]
que são os pontos abaixo a reta $y=2x$.

\vspace{\baselineskip}
\textbf{b)\;}
As curvas de nível de $f$ são compostas pelos pontos onde $f(x, y) = c$
para alguma constante $c$ na imagem de $f$.

Para determinar a imagem de $f$ observamos que $2x-y$ pode assumir
qualquer valor não negativo. Portanto, $\sqrt{2x-y}$ está no intervalo $[0, \infty)$.
Concluímos que os valores de $f$ estão em $[-2, \infty)$, isto é,
\(\Im(f) = [-2, \infty)\)

Para qualquer $c\in\Im(f)=[-2, \infty)$, temos a curva de nível
\begin{align*}
  f(x, y)       & = c              \\
  \sqrt{2x-y}-2 & = c              \\
  \sqrt{2x-y}   & = c + 2          \\
  2x-y          & = (c + 2)^2      \\
  -y            & = (c + 2)^2 - 2x \\
  y             & = 2x - (c + 2)^2
\end{align*}
Que são retas.
Escolhendo os valores $c=-2$, $c=0$ e $c=2$, temos as curvas de nível
\begin{align*}
  \gamma_1\colon & y = 2x - (-2 + 2)^2 = 2x - 0 = 2x \\
  \gamma_2\colon & y = 2x - ( 0 + 2)^2 = 2x - 4      \\
  \gamma_3\colon & y = 2x - ( 2 + 2)^2 = 2x - 16
\end{align*}

\vspace{\baselineskip}
\textbf{c)\;}
Calculando o limite
\setlength{\jot}{12pt}
\setlength{\mathindent}{0cm}
\begin{align*}
  L
  & = \lim_{(x, y)\to(2, 0)} \frac{f(x, y)}{2x-y-4} \\
  & = \lim_{(x, y)\to(2, 0)} \frac{\sqrt{2x-y}-2}{2x-y-4} \\
  & = \lim_{(x, y)\to(2, 0)} \frac{\sqrt{2x-y}-2}{2x-y-4} \times
      \frac{\sqrt{2x-y}+2}{\sqrt{2x-y}+2} \\
  & = \lim_{(x, y)\to(2, 0)} \frac{\left(\sqrt{2x-y}\right)^2-2^2}
                                  {\left(2x-y-4\right)\left(\sqrt{2x-y}+2\right)} \\
  & = \lim_{(x, y)\to(2, 0)} \frac{2x-y-4}{\left(2x-y-4\right)\left(\sqrt{2x-y}+2\right)} \\
  & = \lim_{(x, y)\to(2, 0)} \frac{1}{\sqrt{2x-y}+2} \\
  & = \frac{1}{\sqrt{2\times 2-0}+2} \\
  & = \frac{1}{\sqrt{4} + 2} \\
  & = \frac{1}{4}
\end{align*}
\end{multicols*}
\end{answer}

% 02
%------------------------------------------------------------------------------%
\exercise{25}
Considerando que a equação
\(
  z^3 - xy + yz + y^3 - 2 = 0
\)
define $z$ como função de $x$ e $y$.
\begin{enumerate}[label=\alph*)]
  \item Calcule as derivadas \(\D{z}{x}\) e \(\D{z}{y}\)
  \item Encontre os valores de \(\D{z}{x}\) e \(\D{z}{y}\) em \((1, 1, 1)\)
  \item Construa a aproximação linear da função $z(x, y)$ no ponto \((1, 1)\)
\end{enumerate}
\clearpagequestiononly

\begin{answer}
\vspace{\baselineskip}
\textbf{a)\;}
Considerando $z=z(x, y)$ e derivando por $x$ os dois lados da equação temos
\begin{align*}
  \left(z(x, y)\right)^3 - xy + yz(x, y) + y^3 & = 2 \\
  \D{}{x}\left(\left(z(x, y)\right)^3 - xy + yz(x, y) + y^3\right) & = \D{2}{x} \\
  \D{}{x}\left(\left(z(x, y)\right)^3\right) -
  \D{}{x}\left(xy\right) +
  \D{}{x}\left(yz(x, y)\right) +
  \D{}{x}\left(y^3\right) & = 0 \\
  3\left(z(x, y)\right)^2\D{z}{x} -y +y\D{z}{x} + 0 = 0 \\
  \left(3\left(z(x, y)\right)^2 + y\right)\D{z}{x} = y \\
  \D{z}{x} = \frac{y}{3\left(z(x, y)\right)^2 + y}
\end{align*}
Derivando por $y$ os dois lados da equação temos
\begin{align*}
  \left(z(x, y)\right)^3 - xy + yz(x, y) + y^3 & = 2 \\
  \D{}{y}\left(\left(z(x, y)\right)^3 - xy + yz(x, y) + y^3\right) & = \D{2}{y} \\
  \D{}{y}\left(\left(z(x, y)\right)^3\right) -
  \D{}{y}\left(xy\right) +
  \D{}{y}\left(yz(x, y)\right) +
  \D{}{y}\left(y^3\right) & = 0 \\
  3\left(z(x, y)\right)^2\D{z}{y} -x + \D{y}{y}z(x, y) + y\D{z}{y} + 3y^2 & = 0 \\
  \left(3\left(z(x, y)\right)^2+y\right)\D{z}{y} -x +z(x, y) +3y^2 & = 0 \\
  \left(3\left(z(x, y)\right)^2+y\right)\D{z}{y} & = x -z(x, y) -3y^2 \\
  \D{z}{y} & = \frac{x -z(x, y) -3y^2}{3\left(z(x, y)\right)^2+y}
\end{align*}

\vspace{\baselineskip}
\textbf{b)\;}
Avaliando as derivadas no ponto \((1, 1, 1)\)
\begin{align*}
  \D{z}{x}(1, 1)
  & = \left(\frac{y}{3z^2 + y}\right)\evalat{(1, 1, 1)}{}
    = \frac{1}{3\times 1^2 + 1}
    = \frac{1}{4}
  \\
  \D{z}{y}(1, 1)
  & = \left(\frac{x -z -3y^2}{3z^2+y}\right)\evalat{(1, 1, 1)}{}
    = \frac{1 -1 -3\times 1^2}{3\times 1^2 + 1}
    = \frac{-3}{4}
\end{align*}

\vspace{\baselineskip}
\textbf{c)\;}
A aproximação linear da função $z(x, y)$ no ponto \((x_0, y_0) = (1, 1)\) é
\begin{align*}
  L(x, y)
  & = z(x_0, y_0) + (x-x_0)\D{z}{x}(x_0, y_0) + (y-y_0)\D{z}{y}(x_0, y_0) \\
  & = z(1, 1) + (x-1)\D{z}{x}(1, 1) + (y-1)\D{z}{y}(1, 1) \\
  & = 1 + \frac{1}{4}(x-1) - \frac{3}{4}(y-1)
\end{align*}
\end{answer}

% 03
%------------------------------------------------------------------------------%
\exercise{25}
Encontre os valores máximo e mínimo da função
\(
  f(x, y) = 4xy
\)
restrita a região fechada limitada
\(
  x^2 + y^2 \leq 32
\)
\clearpagequestiononly

\begin{answer}
\setlength\columnsep{3em}
\setlength{\jot}{5pt}
\begin{multicols*}{2}
  \raggedcolumns
  Como a função é contínua e a região é fechada e limitada sabemos que
  $f$ assume um valor máximo e um valor mínimo na região.

  \vspace{\baselineskip}
  Temos que buscar os pontos críticos no interior e aplicar multiplicadores
  de Lagrange na fronteira.

  \vspace{\baselineskip}
  \textbf{Gradiente} de $f$
  \[
    \D{f}{x}
    = \D{}{x}\left(4xy\right)
    = 4y
  \]
  \[
    \D{f}{y}
    = \D{}{y}\left(4xy\right)
    = 4x
  \]

  \vspace{\baselineskip}
  \textbf{Pontos críticos}:
  Impondo \(\nabla f = 0\) temos o ponto interior
  \[
    P_1 = (0, 0)
  \]

  \vspace{\baselineskip}
  \textbf{Multiplicadores de Lagrange}:
  Temos a restrição \(g(x, y) = x^2 + y^2 \leq 9\) e seu gradiente é
  \[
    \D{g}{x}
    = \D{}{x}\left(x^2 + y^2\right)
    = 2x
  \]
  \[
    \D{g}{y}
    = \D{}{y}\left(x^2 + y^2\right)
    = 2y
  \]

  Precisamos resolver o sistema
  \begin{gather*}
    4y = 2\lambda x \\
    4x = 2\lambda y \\
    x^2 + y^2 = 32
  \end{gather*}
  que pode ser simplificado para
  \begin{gather*}
    2y = \lambda x \\
    2x = \lambda y \\
    x^2 + y^2 = 32
  \end{gather*}
  Se $\lambda=0$ temos que $y=0$ pela primeira equação e $x=0$ pela segunda,
  mas esses valores não resolvem a terceira equação. Concluímos que $\lambda\neq0$.

  Assim podemos isolar $x$ na primeira equação e substituir na segunda
  \begin{align*}
    2y                  & = \lambda x          \\
    x                   & = \frac{2y}{\lambda} \\[3mm]
    2x                  & = \lambda y          \\
    2\frac{2y}{\lambda} & = \lambda y          \\
    4y                  & = \lambda^2 y
  \end{align*}
  As soluções são $y=0$ ou $\lambda=\pm 2$.

  Se $y=0$ temos \(4x = 2\lambda y = 0\), que não resolve a terceira equação.

  Se $\lambda=2$, pela primeira equação, temos que $x=y$.
  A terceira equação se torna
  \begin{align*}
    x^2 + y^2 & = 32    \\
    2x^2      & = 32    \\
    x^2       & = 16    \\
    x         & = \pm 4
  \end{align*}
  Encontramos os pontos
  \[
    P_2 = (4, 4)
    \qquad
    P_3 = (-4, -4)
  \]

  Se $\lambda=-2$, pela primeira equação, temos que $x=-y$.
  A terceira equação se torna
  \begin{align*}
    x^2 + y^2    & = 32    \\
    x^2 + (-x)^2 & = 32    \\
    2x^2         & = 32    \\
    x^2          & = 16    \\
    x            & = \pm 4
  \end{align*}
  Encontramos os pontos
  \[
    P_4 = (4, -4)
    \qquad
    P_5 = (-4, 4)
  \]

  \textbf{Avaliando} a função nos pontos encontrados
  \begin{gather*}
    f( 0,  0) = 4\times   0  \times   0  =   0 \\
    f( 4,  4) = 4\times   4  \times   4  =  64 \\
    f(-4, -4) = 4\times (-4) \times (-4) =  64 \\
    f(-4,  4) = 4\times (-4) \times   4  = -64 \\
    f( 4, -4) = 4\times   4  \times (-4) = -64
  \end{gather*}

  O valor mínimo de $f$ é $-64$ e o valor máximo é 64
\end{multicols*}
\end{answer}

%------------------------------------------------------------------------------%
\end{document}
%------------------------------------------------------------------------------%
